\documentclass[UTF8,fontset=fandol,aspectratio=169,11pt]{ctexbeamer}
% Shared Beamer preamble for Big Data and Management Decision-Making slides.

\usetheme{Madrid}
\usecolortheme{default}
\usefonttheme{professionalfonts}

\usepackage{amsmath,amssymb,mathtools}
\usepackage{booktabs,array}
\usepackage{tikz}
\usepackage{graphicx}
\usepackage{listings}
\usepackage{xcolor}
\usetikzlibrary{arrows.meta,positioning,shapes.geometric}

\newcommand{\BDMFandolPath}{/usr/local/texlive/2025/texmf-dist/fonts/opentype/public/fandol/}
\setCJKmainfont[
  Path=\BDMFandolPath,
  UprightFont=FandolSong-Regular.otf,
  BoldFont=FandolSong-Regular.otf,
  ItalicFont=FandolKai-Regular.otf,
  BoldItalicFont=FandolKai-Regular.otf
]{FandolSong-Regular.otf}
\setCJKsansfont[
  Path=\BDMFandolPath,
  UprightFont=FandolSong-Regular.otf,
  BoldFont=FandolSong-Regular.otf
]{FandolSong-Regular.otf}
\setCJKmonofont[
  Path=\BDMFandolPath,
  UprightFont=FandolSong-Regular.otf,
  BoldFont=FandolSong-Regular.otf
]{FandolSong-Regular.otf}
\setmonofont{Menlo}

\definecolor{AcademicBlue}{RGB}{22,44,73}
\definecolor{AcademicTeal}{RGB}{22,119,119}
\definecolor{AcademicGray}{RGB}{76,84,95}
\definecolor{AcademicLight}{RGB}{245,247,250}
\definecolor{AcademicAmber}{RGB}{181,111,35}
\definecolor{CodeBg}{RGB}{248,248,248}

\setbeamercolor{structure}{fg=AcademicBlue}
\setbeamercolor{frametitle}{fg=white,bg=AcademicBlue}
\setbeamercolor{title}{fg=white,bg=AcademicBlue}
\setbeamercolor{block title}{fg=white,bg=AcademicBlue}
\setbeamercolor{block body}{bg=AcademicLight,fg=black}
\setbeamertemplate{navigation symbols}{}
\setbeamertemplate{footline}[frame number]
\setbeamertemplate{itemize item}{\raisebox{0.2ex}{\(\blacktriangleright\)}}
\setbeamerfont{title}{series=\mdseries}
\setbeamerfont{subtitle}{series=\mdseries}
\setbeamerfont{frametitle}{series=\mdseries}
\setbeamerfont{block title}{series=\mdseries}
\setbeamerfont{section in toc}{series=\mdseries}

\lstdefinestyle{pythonstyle}{
  basicstyle=\ttfamily\scriptsize,
  backgroundcolor=\color{CodeBg},
  keywordstyle=\color{AcademicBlue},
  commentstyle=\color{AcademicGray},
  stringstyle=\color{AcademicAmber},
  showstringspaces=false,
  breaklines=true,
  frame=single,
  rulecolor=\color{black!20},
  columns=fullflexible
}

\lstdefinestyle{sqlstyle}{
  language=SQL,
  basicstyle=\ttfamily\scriptsize,
  backgroundcolor=\color{CodeBg},
  keywordstyle=\color{AcademicBlue},
  commentstyle=\color{AcademicGray},
  stringstyle=\color{AcademicAmber},
  showstringspaces=false,
  breaklines=true,
  frame=single,
  rulecolor=\color{black!20},
  columns=fullflexible
}

\lstset{style=pythonstyle}

\hypersetup{
  colorlinks=true,
  linkcolor=AcademicBlue,
  urlcolor=AcademicTeal,
  pdfauthor={大数据与管理决策基础},
  pdfsubject={大数据与管理决策基础课程课件}
}


\title[统计描述与 Dashboard]{第5讲：统计描述、经营波动与 Dashboard}
\subtitle{中心、离散、异常、趋势与管理沟通}
\author{《大数据与管理决策基础》}
\date{}

\begin{document}

\begin{frame}
  \titlepage
  \begin{center}
    \small 核心命题：Dashboard 不是图表堆叠，而是可信指标面向管理行动的组织方式。
  \end{center}
\end{frame}

\begin{frame}{本讲学习目标}
\begin{enumerate}
  \item 解释描述统计在管理决策中的作用；
  \item 区分均值、中位数、标准差、IQR 和变异系数；
  \item 使用移动平均描述经营节奏；
  \item 理解图表是数据到视觉通道的编码；
  \item 区分战略、分析和运营 dashboard；
  \item 生成一个可复核的 dashboard 原型。
\end{enumerate}
\end{frame}

\begin{frame}{课程主线中的第5周}
\begin{center}
\resizebox{0.94\linewidth}{!}{%
\begin{tikzpicture}[
  node/.style={draw=AcademicBlue!65, fill=AcademicLight, rounded corners=1pt,
               minimum height=0.8cm, minimum width=2.2cm, align=center},
  arrow/.style={-{Stealth[length=2mm]}, thick, AcademicTeal}
]
\node[node] (metric) {可信指标};
\node[node, right=0.9cm of metric] (stat) {描述统计};
\node[node, right=0.9cm of stat] (vis) {图表编码};
\node[node, right=0.9cm of vis] (interp) {经营解释};
\node[node, right=0.9cm of interp] (act) {管理行动};
\draw[arrow] (metric) -- (stat);
\draw[arrow] (stat) -- (vis);
\draw[arrow] (vis) -- (interp);
\draw[arrow] (interp) -- (act);
\end{tikzpicture}
}
\end{center}
\end{frame}

\begin{frame}{描述统计回答三个问题}
\begin{table}
\small
\centering
\begin{tabular}{p{0.24\linewidth}p{0.50\linewidth}}
\toprule
问题 & 统计工具\\
\midrule
典型水平是多少？ & 均值、中位数、分位数。\\
波动有多大？ & 方差、标准差、变异系数。\\
是否存在异常？ & IQR、箱线图、目标线、质量提示。\\
\bottomrule
\end{tabular}
\end{table}
\end{frame}

\begin{frame}{中心与离散}
\[
\bar x=\frac{1}{n}\sum_{i=1}^{n}x_i,
\qquad
s^2=\frac{1}{n-1}\sum_{i=1}^{n}(x_i-\bar x)^2.
\]
\[
CV=\frac{s}{\bar x}.
\]

\begin{block}{管理解释}
均值描述典型水平，标准差描述绝对波动，变异系数描述相对波动。
\end{block}
\end{frame}

\begin{frame}{分位数与 IQR}
\[
IQR=Q_3-Q_1.
\]

\begin{block}{IQR 异常规则}
若
\[
x<Q_1-1.5IQR
\quad \text{或}\quad
x>Q_3+1.5IQR,
\]
则将其标记为需要下钻的异常观测。
\end{block}

\begin{alertblock}{注意}
异常识别不是因果解释，只是提出下钻问题。
\end{alertblock}
\end{frame}

\begin{frame}{比率指标必须显示分母}
总体比率：
\[
R=\frac{\sum_g N_g}{\sum_g D_g}.
\]

\begin{table}
\small
\centering
\begin{tabular}{p{0.30\linewidth}p{0.50\linewidth}}
\toprule
Dashboard 显示 & 管理风险\\
\midrule
只显示百分比 & 小样本分组可能被误读。\\
显示百分比和分母 & 能判断比较是否稳定。\\
显示质量状态 & 能判断指标是否可发布。\\
\bottomrule
\end{tabular}
\end{table}
\end{frame}

\begin{frame}{移动平均与经营节奏}
\[
MA_t^{(k)}=\frac{1}{k}\sum_{j=0}^{k-1}x_{t-j}.
\]

\begin{block}{实践规则}
移动平均用于看趋势，但不能替代原始观测；异常值需要同时保留。
\end{block}
\end{frame}

\begin{frame}{第5周样例统计结果}
\begin{table}
\small
\centering
\begin{tabular}{p{0.36\linewidth}p{0.30\linewidth}}
\toprule
指标 & 数值\\
\midrule
总收入 & 71350\\
平均日收入 & 5945.83\\
收入中位数 & 5100\\
收入 IQR & 2875\\
准时交付率 & 0.5833\\
平均延期天数 & 0.75\\
质量问题数 & 6\\
\bottomrule
\end{tabular}
\end{table}
\end{frame}

\begin{frame}{图表是视觉编码}
\begin{table}
\small
\centering
\begin{tabular}{p{0.24\linewidth}p{0.50\linewidth}}
\toprule
组成 & 含义\\
\midrule
data & 表格记录与字段。\\
transform & 聚合、过滤、分箱、移动计算。\\
mark & 点、线、条、文本等图形标记。\\
encoding & 字段到位置、颜色、大小等通道的映射。\\
scale/guide & 坐标轴、图例和尺度。\\
\bottomrule
\end{tabular}
\end{table}
\end{frame}

\begin{frame}{任务与图表选择}
\begin{table}
\small
\centering
\begin{tabular}{p{0.25\linewidth}p{0.46\linewidth}}
\toprule
任务 & 图表\\
\midrule
时间趋势 & 折线图、柱线组合。\\
类别比较 & 排序条形图。\\
分布形态 & 直方图、箱线图。\\
目标达成 & 子弹图、目标线。\\
异常定位 & 散点图、标注点。\\
\bottomrule
\end{tabular}
\end{table}
\end{frame}

\begin{frame}{Dashboard 信息架构}
\[
\text{Overview}
\rightarrow
\text{Compare}
\rightarrow
\text{Explain}
\rightarrow
\text{Drill down}
\rightarrow
\text{Act}.
\]

\begin{block}{最低组成}
KPI 卡片、趋势视图、比较视图、异常视图、质量提示、下钻入口。
\end{block}
\end{frame}

\begin{frame}{Dashboard 类型}
\begin{table}
\small
\centering
\begin{tabular}{p{0.20\linewidth}p{0.52\linewidth}}
\toprule
类型 & 设计重点\\
\midrule
战略型 & 目标达成、趋势、少量高层 KPI。\\
分析型 & 分组比较、下钻、分布和异常解释。\\
运营型 & 近实时状态、任务队列、阈值告警。\\
\bottomrule
\end{tabular}
\end{table}
\end{frame}

\begin{frame}[fragile]{代码入口}
\begin{lstlisting}[language=Python]
from bdm_decision.cases.week05_statistics_dashboard import (
    load_daily_operations,
    descriptive_summary,
    dashboard_kpis,
)

rows = load_daily_operations()
revenue = [row.revenue for row in rows]
print(descriptive_summary(revenue))
print(dashboard_kpis(rows))
\end{lstlisting}
\end{frame}

\begin{frame}{第5周实验交付}
\begin{enumerate}
  \item 计算收入的均值、中位数、标准差、IQR；
  \item 识别至少一个异常日期；
  \item 设计包含 4 个以上 KPI 卡片的 dashboard；
  \item 展示一个趋势视图和一个比较或分布视图；
  \item 写明口径、分母、数据质量状态和下钻路径；
  \item 提交一页管理解释。
\end{enumerate}
\end{frame}

\begin{frame}{本讲小结}
\begin{enumerate}
  \item 描述统计是管理解释的第一层证据；
  \item dashboard 必须显示趋势、比较、分母、异常和质量状态；
  \item 图表设计应尊重视觉感知和任务目标；
  \item 异常不是因果解释，只是下钻入口；
  \item 下一讲进入统计推断与 A/B 测试。
\end{enumerate}
\end{frame}

\end{document}

